% Cross-domain example: computer graphics (primitive-based neural rendering) -> radiative-transfer physics.
% Query 10.48550_arxiv.2603.02887 (SIR-4 CS test, cross-field stratum); gold 10.1364/josaa.18.001929.
% Numbers from colab_qualitative_cs.ipynb (8 Sep 2026; hops_frame_on / hops_frame_off / hops_openie):
%   rank of the gold: Qwen3 cosine 8 -> multi-view scorer 7 -> +OpenIE graph 3 -> +SciAffordGraph 2 -> +CCMP 2;
%   graph channel alone 1; OpenIE graph channel 3,169 (no route within six hops); path weight 16.08 with CCMP, 19.12 without.
% Preamble needs: xcolor, tcolorbox[most], booktabs, float, tikz (positioning, arrows.meta, calc); \resulttablestyle.
\providecolor{hlblue}{RGB}{214,228,247}
\providecolor{hlgreen}{RGB}{214,240,224}
\providecolor{hlorange}{RGB}{252,232,200}
\providecolor{hlpurple}{RGB}{228,220,246}
\providecolor{hlgrey}{RGB}{236,236,238}
\providecommand{\hl}[2]{\colorbox{#1}{\strut #2}}
\providecommand{\gw}[1]{{\color{black!55}(#1)}}
\providecommand{\ty}[1]{{\scriptsize\scshape\bfseries\color{black!75}#1}}
\providecommand{\qbox}[2]{%
  \begin{tcolorbox}[enhanced, colback=gray!4, colframe=black!70, boxrule=0.6pt, arc=2.5pt,
                    left=6pt, right=6pt, top=4pt, bottom=4pt, before skip=4pt, after skip=4pt]
  \textbf{#1}\quad #2
  \end{tcolorbox}}
\providecommand{\ladder}[5]{%
  {\footnotesize rank of this paper:\; Qwen3 cosine \textbf{#1} $\rightarrow$ multi-view scorer \textbf{#2}
  $\rightarrow$ {+}\,OpenIE graph \textbf{#3} $\rightarrow$ {+}\,SciAffordGraph \textbf{#4} $\rightarrow$ {+}\,CCMP \textbf{#5}}}

% =================================================================================================
\begin{figure}[H]
\centering
\small
\qbox{Research problem (Computer Science: computer graphics)}{``Is it possible to replace the
\hl{hlblue}{exponential transmittance function} used in primitive-based neural rendering with alternative
decay models that \hl{hlgreen}{decay faster}, thereby reducing the number of primitives composited per pixel
and accelerating rendering, without compromising image fidelity? [\ldots] Previous optimisation attempts
\hl{hlorange}{assume the exponential model}. Meanwhile, in volumetric rendering of correlated media,
\hl{hlpurple}{non-exponential transmittance functions} have been introduced to account for spatial
correlation, but have not been integrated into discrete primitive compositing.''}

\qbox{Hypothetical answer $h_0$}{``A \hl{hlpurple}{generalized transmittance model employing a polynomial
decay function}, where the opacity is adjusted according to a user-defined curve, allowing for
\hl{hlgreen}{faster decay rates} while preserving colour fidelity through a modified alpha compositing
algorithm.''\quad {\footnotesize (best of eight views; cosine to the gold 0.49, raw query 0.44. Views
4 to 7 are explicit cross-domain analogies: non-linear drag in fluid dynamics, non-exponential decay
envelopes in audio, hyperbolic predation decay in ecology.)}}

\qbox{Retrieved inspiration (Physics: radiative transfer in random media, 2001)}{\emph{On the extinction of radiation by a homogeneous
but spatially correlated random medium} (J.\ Opt.\ Soc.\ Am.\ A). Exponential extinction, the
Beer--Lambert law, follows from Poisson statistics of extinction events in a
\hl{hlorange}{homogeneous medium}; when the medium is \hl{hlpurple}{spatially correlated}, attenuation with
depth is \hl{hlblue}{no longer exponential}.\\[2pt]
\ladder{8}{7}{3}{2}{2}{\footnotesize; graph channel alone \textbf{1}.}}

\caption{Cross-domain example, SIR-4 CS, Computer Science to Physics. A rendering question is answered by radiative-transfer physics. The
query never names Beer--Lambert, Poisson statistics, or random media. The hypothetical answers move the
paper from 8 to 7, the entity graph to 3, and the SciAffordGraph reasoner ranks it first on its own
(fused 2). The route (Figure~\ref{fig:graph-cs-render}, Table~\ref{tab:paths-cs-render}) is the physics
argument itself: the exponential law is a consequence of assuming a homogeneous medium, and the paper's
stochastic treatment of extinction is the method that removes that assumption.}
\label{fig:qual-cs-render}
\end{figure}

% =================================================================================================
\begin{figure}[H]
\centering
\begin{tikzpicture}[
  font=\small, every node/.style={align=left},
  fr/.style={draw=black!70, line width=0.5pt, rounded corners=2.5pt, inner sep=5pt, text width=50mm},
  lim/.style={fr, fill=hlorange}, met/.style={fr, fill=hlpurple}, ent/.style={fr, fill=hlgrey, text width=26mm, inner sep=3pt},
  qry/.style={fr, fill=white, dashed, text width=150mm, align=center},
  gold/.style={fr, fill=hlblue, line width=1.1pt, text width=70mm, align=center},
  e/.style={-{Latex[length=2.2mm]}, line width=0.7pt, color=black!80},
  seed/.style={-{Latex[length=2mm]}, line width=0.6pt, color=black!45, dashed},
  rel/.style={font=\footnotesize\itshape, inner sep=2pt},
  note/.style={font=\scriptsize, color=black!60, text width=34mm},
  panel/.style={font=\small\bfseries, color=black!80},
]
% ---- query
\node[qry] (q) at (0,0) {\ty{query: Computer Science (computer graphics)}\\ ``Replace the exponential transmittance of primitive-based neural rendering with a faster-decaying model, without losing fidelity''};
% ---- left panel: SciAffordGraph
\node[panel] at (-4.2,-1.8) {A. SciAffordGraph route (graph channel rank 1)};
\node[lim] (l1) at (-5.4,-3.1) {\ty{limitation}\\ assumes homogeneous medium};
\node[met] (m1) at (-5.4,-5.7) {\ty{method}\\ stochastic approach to exponential extinction (Poisson statistics of extinction events)};
\node[gold] (g) at (-4.6,-8.7) {\ty{gold paper: Physics (radiative transfer, atmospheric optics)}\\ \emph{On the extinction of radiation by a homogeneous but spatially correlated random medium}};
\draw[seed] (q.south) -- ++(0,-0.35) -| node[pos=0.3, above=3pt, font=\scriptsize, color=black!55] {seed frame of the query} (l1.north);
\draw[e] (l1) -- node[rel, right] {overcomes$^{-1}$} (m1);
\draw[e] (m1) -- node[rel, right] {contributes$^{-1}$} (g);
\node[note, anchor=north west] at ($(l1.north east)+(0.25,0)$) {the query's ``assumes the exponential model'' seeds this frame: the exponential law is what a homogeneous medium implies};
\node[note, anchor=north west] at ($(m1.north east)+(0.25,0)$) {path weight 16.08 with CCMP \gw{19.12 without}; both hops gate 1.00};
% ---- right panel: OpenIE
\node[panel] at (5.4,-1.8) {B. OpenIE entity graph (graph channel rank 3{,}169)};
\node[ent] (e1) at (4.0,-3.1) {neural rendering};
\node[ent] (e2) at (6.9,-3.1) {exponential function};
\node[ent] (e3) at (4.0,-4.3) {alpha blending};
\node[ent] (e4) at (6.9,-4.3) {volume rendering};
\node[ent] (e5) at (4.0,-5.5) {correlated noise};
\node[ent] (e6) at (6.9,-5.5) {reflectance};
\draw[seed] (q.south) -- ++(0,-0.35) -| ($(e1.north)!0.5!(e2.north)$);
\node[gold, text width=60mm, fill=white, draw=black!50, dashed, line width=0.6pt] (g2) at (5.4,-8.7) {\ty{gold paper}\\ \emph{On the extinction of radiation \ldots}\\[1pt] {\footnotesize no route from any seed entity within six hops}};
\draw[-{Latex[length=2mm]}, line width=0.7pt, color=black!50, dotted] ($(e5.south)!0.5!(e6.south)+(0,-0.1)$) -- node[right, font=\scriptsize, color=black!55, pos=0.5] {\;is\_mentioned\_in chains only;\\ \;none reach the paper} (g2.north);
\end{tikzpicture}
\caption{Graph view of the rendering example. \textbf{A.} On the SciAffordGraph the query's
\textsc{limitation} seed ``assumes homogeneous medium'' is overcome by the \textsc{method} the gold
contributes, two typed hops. \textbf{B.} On the OpenIE entity graph built from the same corpus the query's
seed entities are surface terms (``neural rendering'', ``alpha blending'', ``exponential function'') and no
co-mention chain reaches the paper within the reasoner's depth. Dashed arrows are the seeding step; solid
arrows are graph edges; $r^{-1}$ is the inverse relation.}
\label{fig:graph-cs-render}
\end{figure}

% =================================================================================================
\begin{table}[H]
\centering
\resulttablestyle
\setlength{\tabcolsep}{5pt}
\renewcommand{\arraystretch}{1.25}
\begin{tabular}{@{}p{0.11\linewidth} p{0.87\linewidth}@{}}
\toprule
\multicolumn{2}{@{}l}{\textbf{Paths}\quad faster-than-exponential transmittance $\rightarrow$ \emph{On the extinction of radiation by a homogeneous but spatially correlated random medium}} \\
\midrule
SciAfford\-Graph &
16.08 \gw{19.12}:\; (\textsc{limitation}: assumes homogeneous medium,\; \textit{overcomes}$^{-1}$,\;
\textsc{method}: stochastic approach to exponential extinction) $\rightarrow$
(\textsc{method}: stochastic approach to exponential extinction,\; \textit{contributes}$^{-1}$,\; \textbf{gold paper}) \newline
{\footnotesize shortest seed-to-gold distance 2 hops; graph channel rank 1 with and without CCMP; fused rank 2 and 2.} \\
OpenIE &
no route to the gold within six hops (graph channel rank 3{,}169); seeds: neural rendering, exponential function,
alpha blending, volume rendering, correlated noise, reflectance; OpenIE model rank 3 (carried by its scorer channel, rank 3). \\
\bottomrule
\end{tabular}
\caption{Path interpretations for Figure~\ref{fig:qual-cs-render}. Top-weighted route from the query's seed
frames to the gold under the trained reasoner (beam search over the gradient of the graph score with respect
to each layer's edge weights, as in NBFNet and GFM-RAG); $r^{-1}$ is the inverse relation; grey weight is
the same weights with the CCMP gate off. The frame graph states the physical reason in two typed hops. The
entity graph has the paper but no route to it: its edges record which terms a paper mentions, not what
assumption a method removes.}
\label{tab:paths-cs-render}
\end{table}
